\documentclass[12pt,a4paper]{article}
\usepackage[utf8]{inputenc}
\usepackage[T1]{fontenc}
\usepackage[portuguese]{babel}
\usepackage{geometry}
\geometry{a4paper, margin=2.5cm}
\usepackage{hyperref}
\usepackage{graphicx}
\usepackage{xcolor}
\usepackage{listings}

\title{\textbf{Documentação do Simulador de Sistemas Operacionais}\\
\large Guia Detalhado e Didático do Código Fonte}
\author{Desenvolvido para o Projeto 2026}
\date{\today}

\begin{document}

\maketitle
\tableofcontents
\newpage

\section{Introdução ao Projeto}
Este documento explica, de forma clara e acessível, o funcionamento do código-fonte do \textbf{Simulador de Sistemas Operacionais}. O objetivo deste projeto é imitar matematicamente como um computador real toma decisões cruciais nos bastidores, lidando com duas tarefas vitais:
\begin{enumerate}
    \item \textbf{Decidir quem usa o processador (Escalonamento):} Como o cérebro do computador divide sua atenção entre vários programas abertos ao mesmo tempo.
    \item \textbf{Gerenciar o espaço (Memória Virtual):} Como o computador lida com a falta de espaço quando tentamos abrir muitos aplicativos pesados, usando um truque chamado "Paginação" para usar o disco rígido como uma memória de mentira (virtual).
\end{enumerate}

Mesmo que você não tenha experiência prévia com programação, as explicações a seguir usam analogias do dia a dia para que você compreenda como o sistema pensa e age.

\section{Estrutura Geral dos Arquivos}
Um programa de computador geralmente é dividido em várias partes (arquivos), assim como um livro é dividido em capítulos. Neste projeto, a divisão foi feita da seguinte forma:
\begin{itemize}
    \item \texttt{include/} (Os Manuais): Contém arquivos que terminam em \texttt{.h}. Eles são como índices ou resumos que dizem para o sistema "quais ferramentas existem", sem detalhar como elas funcionam por dentro.
    \item \texttt{src/} (As Engrenagens): Contém arquivos que terminam em \texttt{.c} e \texttt{.cpp}. É aqui que o código real (a "mão na massa") está escrito.
    \item \texttt{resources/} (A Matéria Prima): Pastas com os dados que alimentam o sistema, como o arquivo \texttt{processos\_exemplo.csv}, que diz quais programas devem ser simulados.
\end{itemize}

\section{Os Programas: Como eles nascem (\texttt{process.c} e \texttt{process.h})}
Imagine que um processo no computador seja um cliente esperando para ser atendido em um banco. O arquivo \texttt{process.h} define o "crachá" desse cliente, ou seja, as informações que o sistema precisa saber sobre ele.

No código, isso é chamado de uma \textit{estrutura} (Struct), e guarda as seguintes informações:
\begin{itemize}
    \item \textbf{PID (Identificador):} Um número único de senha do cliente.
    \item \textbf{Tempo de Chegada (\textit{Arrival Time}):} A hora exata que o cliente entrou na fila do banco.
    \item \textbf{Tempo de Execução (\textit{Burst Time}):} Quanto tempo o atendimento desse cliente vai demorar no caixa (Processador).
    \item \textbf{Prioridade:} O cliente é um idoso com atendimento preferencial? (Valores menores significam maior prioridade).
    \item \textbf{Memória Necessária:} Quanta "mesa" o cliente precisa para espalhar seus documentos durante o atendimento.
\end{itemize}

O arquivo \texttt{process.c} apenas faz o trabalho braçal de ler o arquivo de planilhas (\texttt{.csv}) e criar os crachás de todos os clientes, colocando-os em uma fila para o sistema atender.

\section{O Caixa do Banco: Escalonamento (\texttt{scheduler.c})}
O arquivo \texttt{scheduler.c} é onde a mágica das filas acontece. Ele funciona como o gerente do banco que decide qual caixa vai atender qual cliente. Como o nosso computador só tem um caixa (um processador único), ele precisa revezar. O sistema implementa três estratégias (algoritmos) diferentes para fazer isso:

\subsection{1. Round-Robin (O Revezamento Justo)}
Nesta estratégia, o gerente diz: \textit{"Todo mundo vai ser atendido por, no máximo, 4 minutos de cada vez"}. Se o seu problema demora 10 minutos, o caixa te atende por 4, pede para você ir para o final da fila, chama o próximo, e você volta depois.
\begin{itemize}
    \item \textbf{Como funciona no código:} O código cria uma fila circular (uma roda gigante). Ele pega o primeiro da fila, roda o relógio até o limite (o \textit{Quantum}). Se o processo terminou, ele é removido. Se não, o código literalmente copia o processo para o final da fila e repete o processo até esvaziar a sala.
\end{itemize}

\subsection{2. SJF Preemptivo (Os Mais Rápidos Primeiro)}
Imagine que você está na fila do mercado com um carrinho cheio, e chega alguém atrás com apenas 1 garrafa de água. O algoritmo \textit{Shortest Job First} (SJF) permite que a pessoa com menos itens passe na frente. E pior (ou melhor), se você já está no caixa passando as compras e a pessoa chega, o caixa para de te atender, guarda suas sacolas e atende a pessoa rápida primeiro.
\begin{itemize}
    \item \textbf{Como funciona no código:} A cada segundo (tempo), o código varre a lista de todos os processos que já chegaram. Ele procura aquele que tem o \textbf{menor tempo restante}. O processo com a tarefa mais curta assume o controle. Isso exige o cálculo constante a cada instante de tempo.
\end{itemize}

\subsection{3. Prioridade Preemptiva (Os Vips)}
Parecido com o anterior, mas o tamanho do trabalho não importa. O que importa é a "Carteirinha Vip".
\begin{itemize}
    \item \textbf{Como funciona no código:} Novamente, a cada milissegundo o sistema verifica a lista. Ele escolhe o processo que tem a melhor prioridade (matematicamente, o menor número). Se um processo super prioritário nascer no meio do nada, ele interrompe o que está sendo executado e rouba a CPU para si.
\end{itemize}

\section{A Mesa de Trabalho: Memória Virtual (\texttt{memory.c})}
A memória RAM do computador é muito rápida, mas muito pequena (como a sua mesa de trabalho). O disco rígido (HD/SSD) é enorme, mas muito lento (como o arquivo no fundo da sala). O \texttt{memory.c} simula a tática de trazer para a mesa apenas os papéis que você precisa usar agora (Páginas), guardando o resto no arquivo.

Quando o sistema precisa de um papel que não está na mesa, acontece um erro chamado \textbf{Page Fault} (Falha de Página). Você precisa pausar tudo, ir até o arquivo, pegar o papel e voltar. O objetivo do sistema é minimizar a quantidade de vezes que você levanta da cadeira. Se a mesa lotar, qual papel você devolve pro arquivo? Temos 3 políticas para isso:

\subsection{A Política FIFO (O Primeiro a Entrar, é o Primeiro a Sair)}
A lógica é simples: o papel que está na mesa há mais tempo é o escolhido para voltar pro arquivo. O código usa uma contagem e a posição sequencial para sempre apontar para o papel mais "velho".

\subsection{A Política LRU (\textit{Least Recently Used})}
O LRU (Menos Recentemente Usado) é mais esperto. Ele pensa: \textit{"Vou jogar fora o papel que faz mais tempo que eu não toco"}.
\begin{itemize}
    \item \textbf{Como o código faz isso:} Cada página na memória guarda a data/hora (o \textit{timestamp}) em que ela foi lida pela última vez. O algoritmo varre todas as páginas e encontra a que tem o tempo mais antigo.
\end{itemize}

\subsection{A Política Ótima (A Bola de Cristal)}
Este é o algoritmo perfeito. Ele é impossível na vida real porque exige prever o futuro. Ele se pergunta: \textit{"Qual desses papéis na mesa eu vou demorar mais para precisar novamente?"}.
\begin{itemize}
    \item \textbf{Como o código faz isso (Nossa Invenção):} Como estamos em uma simulação que nós mesmos construímos, o nosso simulador \textit{conhece} o futuro. O simulador constrói uma fila do que o processo fará. Quando a memória lota, o código vasculha essa fila futurística e busca qual dos itens na memória demorará mais a ser invocado, removendo-o. Este algoritmo recentemente passou por uma pesada otimização no nosso sistema para que rodasse dezenas de milhares de vezes mais rápido usando checagens sequenciais.
\end{itemize}

\section{Os Maestros do Espetáculo (\texttt{main.c} e \texttt{mainwindow.cpp})}
Toda essa lógica não tem utilidade se o humano não a visualizar. 

\begin{itemize}
    \item O \texttt{main.c} é feito para uso em \textbf{Terminal (Tela Preta)}. Ele é focado em rapidez, rodando os cálculos de trás das cortinas sem se preocupar em pintar telas bonitas.
    \item O \texttt{mainwindow.cpp} é o \textbf{Gerenciador Visual (Interface Gráfica)}. Desenvolvido usando a tecnologia \textit{Qt}, ele possui painéis, botões e caixas de seleção. Quando você aperta "Simular", ele conversa com nossos arquivos C, entrega as planilhas e aguarda os cálculos. Ao receber de volta, ele desenha quadrados coloridos na tela que formam o \textit{Gráfico de Linha do Tempo (Gráfico de Ga\section{Mergulho no Código: Detalhamento Exaustivo das Funções}
Esta seção disseca profissionalmente cada arquivo de código (\texttt{.h}, \texttt{.c}, \texttt{.cpp}) e todas as suas funções. Explicaremos qual a responsabilidade exata de cada pedaço de software no sistema.

\subsection{1. O Módulo de Processos (\texttt{process.h} e \texttt{process.c})}
Responsável por interpretar a entrada de dados (CSV) e modelar a estrutura fundamental da simulação, dizendo "quem é" cada processo.

\subsubsection{\texttt{process.h}}
Define os esqueletos das informações (Estruturas de Dados).
\begin{itemize}
    \item \textbf{\texttt{ProcessState} (Enumeração):} Uma lista de etiquetas. Define os estados possíveis de um processo: \texttt{STATE\_NEW} (Novo), \texttt{STATE\_READY} (Pronto), \texttt{STATE\_RUNNING} (Rodando), \texttt{STATE\_WAITING} (Esperando), e \texttt{STATE\_TERMINATED} (Finalizado).
    \item \textbf{\texttt{Process} (Struct):} A espinha dorsal. Uma "caixa" de informações que armazena \texttt{pid}, \texttt{arrival\_time}, \texttt{burst\_time}, \texttt{remaining\_time}, \texttt{priority}, e \texttt{memory\_needed}. Também guarda as métricas finais após a simulação acabar (\texttt{waiting\_time}, \texttt{turnaround\_time}).
\end{itemize}

\subsubsection{\texttt{process.c}}
A fábrica que lê as informações e monta os processos de fato.
\begin{itemize}
    \item \textbf{\texttt{load\_processes\_from\_csv(filepath, procs, max):}} Função de Leitura de Arquivo (\textit{I/O}). Ela abre o arquivo CSV usando a função \texttt{fopen}, ignora o cabeçalho e lê linha por linha usando \texttt{fgets}. A função \texttt{strtok} é utilizada para "fatiar" as informações da linha divididas por vírgulas, e a função \texttt{atoi} converte os textos fatiados em números matemáticos reais. O resultado é despejado no vetor \texttt{procs}.
    \item \textbf{\texttt{print\_processes(procs, n):}} Uma função auxiliar (\textit{Helper}). Usada por nós, desenvolvedores, para imprimir no console uma tabela bonita em texto e verificar se a conversão feita pela função acima ocorreu perfeitamente.
\end{itemize}

\subsection{2. O Módulo Escalonador (\texttt{scheduler.h} e \texttt{scheduler.c})}
O maestro do processador. O arquivo \texttt{scheduler.h} apenas dita as regras e estruturas para guardar resultados (\texttt{TimelineEntry} e \texttt{SchedulerResult}), enquanto o \texttt{scheduler.c} implementa as lógicas profundas.

\subsubsection{\texttt{scheduler.c} - Funções Auxiliares}
\begin{itemize}
    \item \textbf{\texttt{add\_timeline(r, pid, t\_start, t\_end):}} Responsável por "filmar" a execução. Cada vez que um processo entra ou sai do processador, essa função anota quem era o processo (\texttt{pid}) e de que momento até qual momento ele rodou (\texttt{t\_start} até \texttt{t\_end}). Isso será vital para a Interface Gráfica desenhar as cores depois.
    \item \textbf{\texttt{compute\_metrics(procs, n, r):}} Função estritamente matemática. Calcula o tempo de \textit{Turnaround} (Tempo de Conclusão - Tempo de Chegada), o \textit{Waiting Time} (Espera, deduzindo o tempo gasto trabalhando) e o \textit{Response Time} (Resposta do clique à primeira ação). Ela gera as médias estatísticas oficiais da simulação.
\end{itemize}

\subsubsection{\texttt{scheduler.c} - Algoritmos (O Núcleo)}
\begin{itemize}
    \item \textbf{\texttt{sched\_round\_robin(procs, n, quantum, r):}} O Revezamento. Implementa uma Fila Circular (\textit{Queue}). Usa as variáveis numéricas \texttt{head} e \texttt{tail} para navegar num array funcionando como uma roda gigante. Retira um processo, roda por um tempo estritamente limitado à variável \texttt{quantum}, e o devolve à fila final com o comando \texttt{ENQUEUE} caso a tarefa ainda não esteja 100\% concluída.
    \item \textbf{\texttt{sched\_sjf\_preemptive(procs, n, r):}} \textit{Shortest Remaining Time First} (SRTF). Possui um loop infinito enquanto existirem processos não terminados. A cada instante de milissegundo do tempo, ele varre todos os processos (\texttt{best = i}) e seleciona estritamente o de menor \texttt{remaining\_time}. Se no meio do processamento surgir um outro processo mais rápido ainda, ocorre a famosa "preempção" (roubo da CPU).
    \item \textbf{\texttt{sched\_priority\_preemptive(procs, n, r):}} Funciona como um gêmeo quase idêntico ao SJF. Porém, no laço de repetição, a condição de vitória e tomada do processador é baseada em quem tem a menor \texttt{priority} (onde número 1 tem prioridade sobre o número 3).
    \item \textbf{\texttt{run\_scheduler(type, ...):}} O grande "Despachante" (\textit{Dispatcher}). O usuário clica na interface escolhendo um nome, a interface manda um código, e a função realiza um \texttt{switch case} para engatilhar um dos 3 algoritmos explicados acima.
\end{itemize}

\subsection{3. O Módulo de Memória Virtual (\texttt{memory.h} e \texttt{memory.c})}
Este é o módulo mais complexo matematicamente e em engenharia de sistemas do projeto. Ele cuida do armazenamento e descarte de quadros (\textit{Paging}).

\subsubsection{\texttt{memory.c} - Gerenciamento de Quadros (Frames)}
\begin{itemize}
    \item \textbf{\texttt{memory\_init(mem, physical, virtual):}} A função criadora. Aloca os "quadrados físicos" (frames) de memória dinamicamente usando a função \texttt{malloc} de acordo com a conversão de MB para KB da RAM requerida pelo usuário (ex: 256MB). Inicialmente, limpa todos com o valor de \texttt{-1} (que significa "Vazio").
    \item \textbf{\texttt{find\_free\_frame(mem):}} Vasculha o vetor colossal da RAM procurando os ditos espaços com \texttt{-1}. Retorna o índice do espaço limpo ou avisa, devolvendo o próprio \texttt{-1}, que a mesa de trabalho física do computador já lotou e alguém precisará ser ejetado para liberar espaço.
    \item \textbf{\texttt{memory\_free\_process(mem, pt, pid):}} Quando um processo avisa ao Escalonador que encerrou seus trabalhos, esta função varre a tabela do processo. Onde existir uma página na memória dele assinalada com "Válido", a função apaga esse pedaço na RAM, devolvendo os recursos ao computador.
\end{itemize}

\subsubsection{\texttt{memory.c} - Lógicas de Substituição (\textit{Page Replacement})}
Quando \texttt{find\_free\_frame} relata que não há mais espaço, o sistema engatilha uma destas funções candidatas a \textit{vítima}:
\begin{itemize}
    \item \textbf{\texttt{victim\_fifo(mem, pt):}} Algoritmo simples que usa um ponteiro assinalado como \texttt{victim\_hint}. Ele avança o ponteiro de forma circular de 1 em 1 para garantir que a página expulsa, matematicamente, seja a mais "antiga" na fila desde que a simulação rodou.
    \item \textbf{\texttt{victim\_lru(mem, pt):}} Na nossa implementação para varreduras simples, usa o conceito temporal e posicional para identificar a página "menos tocada recentemente", funcionando ativamente em conjunto com o registro de acesso.
    \item \textbf{\texttt{victim\_optimal(pt, future\_refs, future\_len, current\_pos):}} O milagre Otimizado da previsão. Esta função varre um vetor profético (\texttt{future\_refs}) usando marcadores para enxergar de antemão quais páginas o aplicativo exigirá nos próximos segundos. Ele intercepta e ejeta a página \texttt{last\_page} com cirúrgica precisão, que será a folha que tardará mais a ser necessária. 
\end{itemize}

\subsubsection{\texttt{memory.c} - Simulação de Acesso e Falhas (\textit{Page Faults})}
\begin{itemize}
    \item \textbf{\texttt{memory\_access\_page(...):}} O guarda de segurança. Toda requisição passa por ele. Tenta ler um pedaço da memória; se já o possuir (\texttt{valid == 1}), apenas sinaliza sucesso (\textit{Page Hit}). Se não tiver (\textit{Page Fault}), procura espaço e rouba a vaga de quem for selecionado como vítima.
    \item \textbf{\texttt{memory\_load\_process(...):}} O Gerador de Stress. Emulações cruas apenas alocam tudo 1 vez. Esta função constrói um array de requisições de 3 vezes o tamanho do processo com fatias randomizadas (por \texttt{rand()}) imitando a desorganização que ocorre com os programas reais na RAM do seu Windows ou Mac, causando falhas de páginas reais.
\end{itemize}

\subsection{4. Interação Humano-Máquina (\texttt{main.c} e \texttt{mainwindow.cpp})}
A alma lógica existe, mas é preciso dar-lhe um rosto. Dividimos isso em duas vertentes.

\subsubsection{\texttt{main.c} (Interface Terminal de Desenvolvedor)}
O arquivo mestre compilado com a instrução de atalho \texttt{make cli}. Ele empacota toda as invocações das bibliotecas descritas acima e formata as saídas estatísticas na "Tela Preta" com a diretiva \texttt{printf}. Sendo livre do gigantesco peso de interfaces gráficas, ele serve ativamente para rodarmos o simulador com alto desempenho para checagem rápida de depuração em \textit{Debugging} de código C.

\subsubsection{\texttt{mainwindow.cpp} (O Frontend Gráfico em Qt C++)}
O rosto visual deslumbrante e moderno escrito no paradigma Orientado a Objetos (C++ e Biblioteca Qt).
\begin{itemize}
    \item \textbf{\texttt{TimelineWidget::paintEvent()}:} A classe artística do simulador de gantt. Usa a biblioteca \texttt{QPainter} para ler os registros numéricos guardados pela \texttt{add\_timeline()}. Ela preenche os pixels da tela convertendo números de tempo em posições X no monitor, pintando os blocos \texttt{fillRect} com uma paleta de 8 cores para diferenciar cada processo.
    \item \textbf{\texttt{MainWindow::buildUI()}:} O construtor civil do projeto. Invoca o \textit{Toolkit UI} arranjando \texttt{QVBoxLayout} e \texttt{QGridLayout} no grid. Insere interações de mouse em \texttt{QPushButton} e numéricas em caixas giratórias \texttt{QSpinBox}.
    \item \textbf{\texttt{MainWindow::onSimulate()}:} O sinal de Ignição. Ativada no evento \textit{click()} do botão, recolhe todos as diretivas gráficas pedidas e orquestra uma viagem de ida e volta para as lógicas matemáticas do módulo em C puro. Pega as \textit{Page Faults} e as injeta de volta nas \texttt{QLabel} de resultados, gerando a interface maravilhosa ao qual estamos acostumados.
\end{itemize}

\section{Conclusão Definitiva}
Com todos estes detalhes e pormenores, compreendemos perfeitamente a profundidade estrutural deste Simulador de Sistemas. A qualidade de todo este maquinário computacional baseia-se em um pilar inquebrável da Engenharia de Software: \textbf{o Baixo Acoplamento}. As equações estatísticas de Memória e Escalonamento estão cegas aos detalhes visuais, guardadas confiavelmente em C. O Simulador visual apenas empurra dados e recebe respostas prontas. O formato modular propicia que nosso simulador performe cálculos tridimensionais avançados na mesma facilidade que rodaria um cálculo de um programa escolar, resultando em um sistema maduro, incrivelmente rápido e estável, preparado para emular as raízes de qualquer Sistema Operacional da vida real.
\end{document}